\documentclass[12pt]{article}
\usepackage[a4paper, margin=2.5cm]{geometry}
\usepackage{setspace}
\usepackage{hyperref}
\usepackage{listings}
\usepackage{xcolor}

\definecolor{codegray}{rgb}{0.5,0.5,0.5}
\definecolor{codepurple}{rgb}{0.58,0,0.82}

\lstdefinestyle{commitstyle}{
    basicstyle=\ttfamily\small,
    breaklines=true,
    frame=single,
    numbers=left,
    numberstyle=\tiny\color{codegray},
    keywordstyle=\color{codepurple},
    showstringspaces=false
}

\onehalfspacing

\begin{document}

\begin{center}
    {\LARGE \textbf{Historial de Commits del Proyecto}}\\[0.5cm]
    {\large Documentación del Control de Versiones}\\[0.2cm]
    {\large Autor: Andrés García}\\
    {\large Fecha: \today}
\end{center}

\vspace{1cm}

\section{Introducción}
Este documento presenta el historial completo de commits del proyecto.
El control de versiones se ha gestionado mediante Git y GitHub, lo que ha
permitido un desarrollo estructurado, la trazabilidad de los cambios y
una evolución progresiva del trabajo realizado.

\section{Proceso de Desarrollo}
Cada commit representa un paso dentro del desarrollo del
proyecto, incluyendo tareas como la preparación de los datos, la corrección
de errores, la implementación de nuevas funcionalidades y la optimización
del código.

El uso de commits frecuentes y descriptivos ha facilitado el seguimiento
del progreso del proyecto y ha permitido identificar de forma clara las
modificaciones realizadas en cada etapa.

\section{Registro de Commits}
A continuación mostramos el listado completo de commits del repositorio,
ordenados cronológicamente desde el inicio del proyecto hasta su estado
actual.

\lstset{style=commitstyle}
\lstinputlisting{commits.txt}

\section{Conclusión}
Como se puede observar en el historial de commits, el proyecto se ha realizado de una manera organizada y estructurada.
El uso de Git y GitHub ha sido fundamental para mantener un control riguroso sobre los cambios realizados, facilitando la colaboración y la gestión del código fuente.

\end{document}
